\section{Nutzung künstlicher Intelligenz im Projekt}

Im Rahmen des Projekts wurden verschiedene KI-Modelle unterstützend eingesetzt. Ihr Einsatz beschränkte sich dabei vor allem auf klar abgegrenzte Aufgaben, bei denen bereits bestehende Konzepte umgesetzt, überprüft oder sprachlich aufbereitet werden mussten.

Bei der Erstellung dieser Dokumentation wurde künstliche Intelligenz insbesondere zur sprachlichen Überarbeitung und Professionalisierung eigener Textentwürfe verwendet. Darüber hinaus unterstützten KI-Modelle bei der Überführung der Inhalte in geeigneten LaTeX-Code, beispielsweise bei der Strukturierung von Abschnitten, der Formatierung von Abbildungen und Tabellen sowie bei der korrekten Verwendung von Referenzen. Die fachlichen Inhalte und die zugrunde liegenden Aussagen wurden jedoch durch das Projektteam erarbeitet und anschließend mithilfe der KI sprachlich und formal verfeinert.

Ein weiterer Anwendungsbereich war die Erstellung und Überarbeitung von Skripten zur Datenaufbereitung. Aufgrund der großen zu verarbeitenden Datenmengen und der notwendigen Transformationen konnten KI-Modelle insbesondere bei der Implementierung wiederkehrender Verarbeitungsschritte, bei der Fehleranalyse sowie bei der Optimierung bestehender Skripte unterstützen. Die Anforderungen an die Daten, die Transformationsregeln und die Überprüfung der erzeugten Ergebnisse wurden weiterhin eigenständig durch das Team festgelegt und durchgeführt.

Auch bei der Entwicklung des Frontends und des Backends kam künstliche Intelligenz punktuell zum Einsatz. Sie unterstützte unter anderem bei der Implementierung bekannter Programmiermuster, der Erstellung einzelner Codeabschnitte, der Fehlersuche und der Verbesserung bereits vorhandener Lösungen. Generierter oder vorgeschlagener Code wurde nicht ungeprüft übernommen, sondern durch das Team nachvollzogen, an die Projektanforderungen angepasst und anschließend getestet.

Die grundlegende Konzeption des Systems, die Modellierung der Daten, die Auswahl der verwendeten Technologien sowie sämtliche Architektur- und Designentscheidungen wurden eigenständig durch das Projektteam getroffen. Gleiches gilt für die Analyse der Anforderungen und die Entwicklung der fachlichen Lösungsansätze. Die KI diente somit hauptsächlich als unterstützendes Werkzeug bei der konkreten Umsetzung bereits erarbeiteter Lösungen und übernahm nicht die eigentliche konzeptionelle oder fachliche Denkarbeit.

Durch den gezielten Einsatz künstlicher Intelligenz konnten insbesondere zeitaufwendige Tätigkeiten wie Formatierung, sprachliche Überarbeitung, Implementierung wiederkehrender Strukturen und technische Fehleranalyse effizienter gestaltet werden. Die Verantwortung über die inhaltliche und technische Korrektheit sowie für die abschließende Bewertung und Integration aller Ergebnisse verblieb dabei vollständig beim Projektteam.